% 本文件由 LaTeX 排版 Agent 根据有效 translation.json 自动生成。
% author=Hippolytus; work_id=0516; title=论基督与敌基督

\chapter{论基督与敌基督}

% structure: p0001 source=h1 target=omitted reason=page-title-already-rendered-by-parent

1. 正如你所愿，我亲爱的弟兄\Person{德奥斐禄}，要全面了解我先前向你概括的那些主题，我认为应当将这些探究之事清晰地摆在你面前，大量取用圣经本身，如同从神圣的泉源中汲取，好使你不仅因人的见证而乐于听闻，还能借由在天主权威的光照下审视它们，在一切事上光荣天主。因这将是我们为你在此生旅途所供应的一种确实的储备，使你藉着巧妙的论证，将大多数人难以理解且领会不当的事，像播在肥沃洁净的土壤里一样，播在你心田的地里。藉着这些，你也能使那些反对并诋毁救恩之道的人缄默。只是你要小心，不可将这些事交给不信和亵渎的舌头，因为那是寻常的危险。但要交给虔诚忠信的人，他们愿意在敬畏中圣洁公义地生活。因为蒙福的宗徒并非无目的地劝勉弟茂德，说：\BibleQuote{弟茂德啊，你要保守所托付你的，躲避世俗的虚谈和那敌真道、似是而非的学问；有人自称有这学问，就在信德上偏离了。} 又说：\BibleQuote{所以，我儿，你要在基督耶稣的恩宠上刚强起来。你在许多劝勉中从我听见的，要交托那忠心能教导别人的人。} 那么，如果蒙福的（宗徒）以虔诚的谨慎传授了这些众人容易知道的事，因为他藉圣神觉察到\BibleQuote{众人不是都有信德}，\BibleRef{得后 3:2}，那么，我们若轻率而不加思索地将天主的启示交给亵渎和不配的人，又将是何等大的危险呢？\par

2. 因为蒙福的先知们，可以说是成了我们的眼睛，他们藉着信德预见了圣言的奥秘，并成为这些事对后代的执行者，不仅报告过去，也宣告现在和将来，好使先知不只为当时，也为万世预言未来，而被视为（真）先知。因为这些先祖被圣神所装备，并大受圣言本身的尊崇；正如乐器一般，他们常有圣言如同拨子，与他们结合，当被圣言感动时，先知们就宣布天主所愿的。因为他们不是凭自己的能力说话（这一点不可误会），也不是宣告自己所喜悦的。而是首先，他们被圣言赋以智慧，然后又藉着异象正确受教于未来。然后，当他们自己如此完全信服时，他们就说出那些仅被天主启示给他们、向所有其他人隐藏的事。因为先知凭什么理由被称为先知，除非他在圣神中预见未来？因为如果先知说任何偶然的事，那么他在说众人眼下之事时就不是先知了。但那人详细阐明尚未发生的事，才被正当地视为先知。因此，先知们从一开始就被正当地称为\Quote{先见者}。\BibleRef{撒上 9:9} 因此，我们也正确受教于他们从前所宣告的，不是凭自己的能力说话。因为我们并不试图在他们从前所说的话中做任何改动，而是将这些话写成的圣经公之于众，并读给那些能正确相信的人听；因为这对双方都有益处：对说话者而言，是正确记住并陈明古时所说的话；对听者而言，是留心所听的话。既然在这项工作中，双方都有任务，即说话者要忠实地、不顾风险地宣告，听者要凭信德听见并接受所说的话，我恳求你与我一同在祈祷中向天主努力。\par

3. 那么，你愿意知道天主的圣言——祂又是天主子，正如祂从太初就是圣言——在古时是如何向蒙福的先知们传达祂的启示的吗？好吧，正如圣言藉着众圣徒显明祂的慈悲和祂毫无偏私，祂光照他们，并使他们适应于对我们有益的事，如同一位高明的医生，了解人的软弱。祂喜爱教导无知的人，使迷途的人转回祂自己的真道。而那些凭信德生活的人，容易寻见祂；对于心地纯洁、圣洁、渴望敲门的人，祂立刻开门。因为祂并不将任何仆人视为不配接受神圣奥秘。祂不因富有而高看富人，也不因贫穷而轻视穷人。祂不藐视外邦人，也不将阉人视为非人。祂不因女人起初的悖逆而憎恨女性，也不因男人的过犯而拒绝男性。但祂寻找一切人，愿意拯救一切人，愿意使众人成为天主的子女，并召叫众圣徒达到一个成全的人。因为只有一个天主子（或仆人），藉着祂，我们接受因圣神而有的重生，也愿意都达到一个成全的、属天的人。\BibleRef{弗 4:13}\par

4. 因为圣言原是无形无体的，却藉着童贞玛利亚取了圣洁的肉身，为自己预备了一件袍子，好像新郎一样，在十字架的苦难中编织而成，为的是藉着将自己的德能与我们有死的身體结合，使不朽坏的与可朽坏的相掺，强壮的与软弱的相融，从而拯救那即将灭亡的人。因此，织机的梁木就是主在十字架上的受难，其上的经线是圣神的力量，纬线是由圣神所编织的圣洁肉身，线是那藉着基督的爱将两者结合并合而为一的恩宠，综框（或杆）是圣言；工人是那些为基督编织这件美丽、长久、完美的长袍的圣祖和先知；圣言像综框（或杆）一样穿过他们，藉着他们完成圣父所愿意的事。\par

5. 但现在时间紧迫，需要立即考虑眼前的问题，而导言中关于天主荣耀的论述已经足够，因此我们应当亲自拿起圣经，从中找出敌基督的来临是什么、是怎样的；在什么时机和什么时间那个可憎者要被揭示；他来自何处、来自哪个支派；他的名字是什么——就是圣经中以数字所指明的；他怎样在人民中行欺骗，从地极聚集他们；他怎样激起对圣徒的迫害和磨难；他怎样自称为神；他的结局是什么；主的突然显现怎样从天上被揭示；全世界的焚毁是什么；圣徒的荣耀和天国是什么样的——当他们与基督一同为王时；以及恶人在火中的惩罚是什么。\par

6. 如今，我们的主耶稣基督（祂也是神）曾以狮子的形象被预言，是因祂的君王身份和荣耀；同样，圣经也早已论及敌基督如同狮子，是因他的暴虐和强横。因为那欺骗者企图在一切事上效法天主子。基督是狮子，所以敌基督也是狮子；基督是君王，\BibleRef{若 18:37}所以敌基督也是君王。救主曾显现为羔羊，\BibleRef{若 1:29}同样，他也要以羔羊的样子出现，但内心却是豺狼。救主在割礼中来到世界，他也同样会来。主派遣宗徒到万民中，他也同样会派遣假宗徒。救主聚集了四散的羊，他也同样会聚集四散的人民。主赐给信祂的人一个印记，他也同样会赐下一个印记。救主以人的形状显现，他也将以人的形状来临。救主复活并显示祂的圣体好像圣殿，\BibleRef{若 2:19}他也将在耶路撒冷建造一座石头的殿。他的欺骗手段我们将在后面展示。但现在让我们转向眼前的问题。\par

7. 蒙福的雅各伯在他的祝福中预言性地论及我们的主和救主，说：\Quote{犹大！你的兄弟要赞美你；你的手必压在仇敌的颈上；你父亲的儿子要向你俯伏。犹大是只幼狮；我儿，你从嫩芽上去了；他蹲伏，他卧如狮子，又如幼狮；谁能唤醒他？权杖必不离犹大，立法者必不离他两腿之间，直到那所应许者来到；他将是万民的期待。他将自己的驴拴在葡萄树上，将自己的驴驹拴在葡萄藤上；他在酒中洗自己的衣服，在葡萄的血中洗自己的袍子。他的眼睛因酒而喜乐，他的牙齿比奶更白。}\par

8. 我知道如何详细解释这些事，但在目前我以为引用这些话语本身是合适的。然而，既然这些话语本身催促我们去谈论它们，我就不省略了。因为这些确实是神圣而荣耀的事，是极有益于灵魂的事。先知使用\Emph{幼狮}一词，是指那按肉身从犹大和达味而出的一位，祂并非由达味的种子所生，而是由圣神的能力感孕，从圣洁的地上嫩芽中出来。因为依撒意亚说：\Quote{由叶瑟的根茎将发出一个嫩枝，由它的根上将开出一朵花。}\BibleRef{依 11:1}依撒意亚称为\Emph{花}的，雅各伯称为嫩芽。因为祂先发出嫩芽，然后在世上开花。而\Quote{他蹲伏，他卧如狮子，又如幼狮}这句话，是指基督三天的安眠（死亡，卧下）；正如依撒意亚也说：\Quote{忠信的熙雍怎么成了妓女！它曾充满正义，其中住宿（卧）着公义；但现在却住着凶手。}\BibleRef{依 1:21}达味也说过同样的话：\Quote{我躺下（卧）安睡，我又醒来，因为上主扶持我。}在这句话中，他指出了他的安眠和复活。雅各伯说：\Quote{谁能唤醒他？}这正是达味和保禄所提及的，正如保禄所说：\Quote{天主父，就是使祂从死者中复活的那一位。}\BibleRef{迦 1:1}\par

9. 而在说到\Quote{权杖必不离犹大，立法者必不离他两腿之间，直到那所应许者来到；他将是万民的期待}时，他将这预言的应验归于基督。因为祂是我们的期待。因为我们期待祂，凭信德我们看见祂带着权能从天而降。\par

10. \Quote{将自己的驴拴在葡萄树上：}这意味着祂将受割礼的人民与祂自己的召叫联合起来。因为祂是葡萄树。\BibleRef{若 15:1}\Quote{将自己的驴驹拴在葡萄藤上：}这表示外邦人的子民，因为祂召叫受割礼的和未受割礼的归于同一信德。\par

11. \Quote{祂要在酒中洗他的衣服，}这是指祂的父藉圣神在约旦河所发的声音。\Quote{在葡萄的血中洗他的长袍。}那么，是什么葡萄的血呢？岂不正是祂自己的血肉，像一串葡萄挂在树上？从祂的肋旁流出了两股泉源：血和水，万民在其中被洗濯和洁净，祂以这些万民为袍子围绕自己。\par

12. \Quote{祂的眼睛因酒而欢喜。}基督的眼睛是什么呢？不就是那些蒙福的先知吗？他们藉圣神预见并预先宣告了将临于祂的苦难，并以属灵的眼睛看见祂在权能中而喜乐，他们由圣言自身和祂的恩宠装备（蒙召）。\par

13. 当他说：\Quote{祂的牙齿比奶更白，}他指的是从基督圣口发出的诫命，这些诫命像奶一样纯洁（洁净）。\par

14. 圣经先前就是这样预表这狮子和幼狮的。同样，我们发现关于假基督也有记载。因为梅瑟这样说：\Quote{丹是幼狮，他要从巴商跃出。}\BibleRef{申 33:22}但为使无人误解以为这是论及救主，请仔细留心此事。他说：\Quote{丹是幼狮；}他提到丹支派，清楚宣告了假基督将出自哪个支派。因为正如基督出自犹大支派，同样假基督将出自丹支派。此事确然，我们从雅各伯的话也看到：\Quote{愿丹成为一条蛇，躺在地上，咬伤马蹄。}\BibleRef{创 49:17}那么，蛇是指什么呢？岂不正是假基督，那欺骗者，在《创世纪》\BibleRef{创 3:1}中提到的，他欺骗了厄娃，并绊倒了亚当（\OriginalTerm{脚跟}{\GreekText{πτερνι}}）？\par

\SourceRecord{Translated by J.H. MacMahon.；Ante-Nicene Fathers；Vol. 5.；Edited by Alexander Roberts, James Donaldson, and A. Cleveland Coxe.；Buffalo, NY: Christian Literature Publishing Co.,；1886.；<http://www.newadvent.org/fathers/0516.htm>.}
